%% Agent 11. Adversarial attack--heal on Stage 1 of the Vol III functor
%% Phi_d = Sp^{ch}_{Sigma_{d-1}, C} o Phi^{FA}_d. Target: the
%% construction of Phi^{FA}_d: CY_d-Cat --> E_d-HolFA via
%% Kontsevich--Tamarkin E_d-formality + Costello--Gwilliam--Li locality.
%% For d=3, the canonical E_3-holomorphic FA is produced by 6D holomorphic
%% Chern--Simons on a CY_3 background coupled to the chosen CY_3 category.
%% Five attack-heal cycles, CG voice, first-principles. Author: Raeez Lorgat.

\documentclass[11pt]{article}
\usepackage[localtheorems]{raeez-math-template}
\newtheorem{theorem}{Theorem}
\newtheorem{proposition}[theorem]{Proposition}
\newtheorem{lemma}[theorem]{Lemma}
\newtheorem{corollary}[theorem]{Corollary}
\newtheorem{scholium}[theorem]{Scholium}
\newtheorem{construction}[theorem]{Construction}
\newtheorem*{attack}{Attack}
\newtheorem*{heal}{Heal}
\newtheorem*{counterexample}{Counterexample}
\theoremstyle{definition}
\newtheorem{definition}[theorem]{Definition}
\newtheorem{remark}[theorem]{Remark}
\newtheorem{example}[theorem]{Example}

\providecommand{\CC}{\mathbb{C}}
\providecommand{\RR}{\mathbb{R}}
\providecommand{\ZZ}{\mathbb{Z}}
\providecommand{\QQ}{\mathbb{Q}}
\providecommand{\PP}{\mathbb{P}}
\providecommand{\HH}{\mathrm{HH}}
\providecommand{\fg}{\mathfrak{g}}
\providecommand{\hbarb}{\hbar}
\providecommand{\cA}{\mathcal{A}}
\providecommand{\cB}{\mathcal{B}}
\providecommand{\cE}{\mathcal{E}}
\providecommand{\cF}{\mathcal{F}}
\providecommand{\cC}{\mathcal{C}}
\providecommand{\cM}{\mathcal{M}}
\providecommand{\cO}{\mathcal{O}}
\providecommand{\cP}{\mathcal{P}}
\providecommand{\cU}{\mathcal{U}}
\providecommand{\cV}{\mathcal{V}}
\providecommand{\Dolb}{\mathrm{Dolb}}
\providecommand{\Disj}{\mathrm{Disj}}
\providecommand{\Obs}{\mathrm{Obs}}
\providecommand{\Opens}{\mathrm{Opens}}
\providecommand{\Sym}{\mathrm{Sym}}
\providecommand{\CE}{\mathrm{CE}}
\providecommand{\Lie}{\mathrm{Lie}}
\providecommand{\Ass}{\mathrm{Ass}}
\providecommand{\Com}{\mathrm{Com}}
\providecommand{\Gerst}{\mathrm{Gerst}}
\providecommand{\hCS}{\mathrm{hCS}}
\providecommand{\Conf}{\mathrm{Conf}}
\providecommand{\Kos}{\mathsf{K}}
\providecommand{\KT}{\mathrm{KT}}
\providecommand{\CGL}{\mathrm{CGL}}
\providecommand{\BM}{\mathrm{BM}}
\providecommand{\BV}{\mathrm{BV}}
\providecommand{\BRST}{\mathrm{BRST}}
\providecommand{\FA}{\mathrm{FA}}
\providecommand{\HolFA}{\mathrm{HolFA}}
\providecommand{\Tot}{\mathrm{Tot}}
\providecommand{\Coh}{\mathrm{Coh}}
\providecommand{\Fuk}{\mathrm{Fuk}}
\providecommand{\MF}{\mathrm{MF}}
\providecommand{\CoHA}{\mathrm{CoHA}}
\providecommand{\opp}{\mathrm{op}}

\title{Stage 1 of $\Phi_d$: $\Phi^{\FA}_d\colon \mathrm{CY}_d\text{-Cat}\to
       \mathbb{E}_d\text{-}\HolFA$ via Kontsevich--Tamarkin + Costello--Gwilliam--Li.
       \\ Six-dimensional holomorphic Chern--Simons as the $d=3$ canonical avatar.
       \\ Five adversarial attack--heal cycles}
\author{Raeez Lorgat}
\date{2026-04-24}
\begin{document}
\maketitle

\begin{abstract}
The Vol III functor
$\Phi_d\colon \mathrm{CY}_d\text{-Cat}^{\mathrm{cyc},\mathrm{prop}}
\to \mathbb{E}_n\text{-ChirAlg}(\cM_d)$
factors through a two-stage composition
\[
\Phi_d \;=\; \mathrm{Sp}^{\mathrm{ch}}_{\Sigma_{d-1}, C}\circ \Phi^{\FA}_d,
\qquad
\Phi^{\FA}_d\colon \mathrm{CY}_d\text{-Cat}^{\mathrm{cyc},\mathrm{prop}}
\longrightarrow \mathbb{E}_d\text{-}\HolFA(X),
\]
with Stage 1 an $\mathbb{E}_d$-holomorphic factorisation algebra on the
Calabi--Yau background $X$ and Stage 2 a specialisation obtained by
factorisation homology over a transverse $(d-1)$-cycle $\Sigma_{d-1}
\subset X$ followed by restriction to a reference curve $C\subset X$.
This paper inscribes Stage 1 as a theorem, canonical up to contractible
choice, via a two-ingredient anchor: \emph{(KT)} Kontsevich--Tamarkin
$\mathbb{E}_d$-formality for the Hochschild cochains of a
$\mathrm{CY}_d$-$A_\infty$-category $\cC$, and \emph{(CGL)} the
Costello--Gwilliam--Li locality theorem for holomorphic factorisation
algebras of observables of a BV-quantised holomorphic theory on $X$.
For $d=3$ the canonical avatar is six-dimensional holomorphic
Chern--Simons on $X$ coupled to $\cC$-matter: Costello's BV
quantisation on $(\CC^3, \Omega_0)$ produces the flat-space
$\mathbb{E}_3$-holomorphic FA (Gwilliam--Williams~2021 gives the
$\mathbb{E}_3$-algebra of observables; Costello--Li~2016 extends by
background-field coupling). Five adversarial cycles stress every load-bearing
step of the construction against primary literature, verify the
dimensional ladder, exhibit the coupling to $\cC$, and pin the
$\mathbb{E}_d$-structure to Dunn--Lurie additivity on the Dolbeault
complex.
\end{abstract}

\section*{Prolegomenon: what Stage 1 is and is not}

Stage 1 is a construction of the form
\begin{equation}
  \cC \;\longmapsto\; \Obs^{\mathrm{BV}}_{\hCS(d)\,\text{coupled to }\cC}
  \;=:\; \Phi^{\FA}_d(\cC)
  \label{eq:stage1-definition}
\end{equation}
where $\hCS(d)$ is the $(2d)$-real-dimensional holomorphic
Chern--Simons theory on a CY$_d$ variety $X$ of shift $d-4$, whose BV
field $\cA \in \Omega^{0,\bullet}(X,\fg)[1]$ is matter-coupled to the
cyclic $A_\infty$-category $\cC$ (the gauge Lie algebra $\fg$ is the
derived global sections of endomorphisms of a generator of $\cC$).
The output is an $\mathbb{E}_d$-holomorphic factorisation algebra on
$X$ in the sense of Costello--Gwilliam \emph{Factorization Algebras in
Quantum Field Theory} Vol II~\S10--11 (hereafter CGII). The
construction is well-defined (canonical up to contractible choice of
renormalisation and propagator gauge) modulo two inputs that feed it
externally:
\begin{itemize}
  \item[(KT)] Kontsevich~1999 (case $d=1$) and Tamarkin~2003 (higher
        $d$): the Hochschild cochain complex $\HH^\bullet(\cC,\cC)$,
        equipped with its natural Gerstenhaber~/~brace $\mathbb{E}_d$-algebra
        structure, is \emph{formal}: there is a zig-zag of
        $\mathbb{E}_d$-algebra quasi-isomorphisms between it and its
        cohomology $H^\bullet(\HH^\bullet(\cC,\cC))$ as an
        $\mathbb{E}_d$-algebra. Formality is proved with a transcendental
        constant from a Drinfeld associator (rational for Kontsevich on
        a formal disc; Drinfeld even associator for Tamarkin in general
        $d$). The choice of associator parametrises a torsor under
        $\mathrm{GRT}_1(\QQ)$.
  \item[(CGL)] Costello~2011 \emph{Renormalization and Effective Field
        Theory} Theorem 13.4.1 (effective BV quantisation of a gauge
        theory on a manifold with nilpotent classical interaction) and
        CGII Theorem~5.0.1 (the factorisation algebra of quantum
        observables carries an $\mathbb{E}_n$-structure under the Dunn
        hypothesis of holomorphic translation invariance); together with
        Li--Li~2020 (the Chern--Simons avatar for CY$_d$ specifically is
        well-defined and carries the $\mathbb{E}_d$-structure after
        Dolbeault-cohomology reduction).
\end{itemize}
Stage 2, which we do \emph{not} construct here, is the specialisation
functor $\mathrm{Sp}^{\mathrm{ch}}_{\Sigma_{d-1},C}$: factorisation
homology of $\Phi^{\FA}_d(\cC)$ over a transverse framed oriented
$(d-1)$-cycle $\Sigma_{d-1}\subset X$, followed by restriction of the
resulting $\mathbb{E}_1$-FA-valued sheaf to a reference curve $C\subset X$.
Stage 2 is an $\mathbb{E}_1$-chiral algebra on $C$; at $d=3$ the
$(\Sigma_2, C) = (K3, E)$ slot produces the $\mathbb{E}_1$-chiral
family on the elliptic curve $E$ that Vol III calls the
$\mathrm{CY}_3$-family (lattice VOA for $\Coh(E)$; Mukai--Heisenberg
for $D^b(K3)$; $Y^+(\widehat{\mathfrak{gl}}_1)$ for $\CoHA(\CC^3)$).

The deliverable is: \emph{Stage 1 is a theorem.} Stage 2 is a theorem
too but elsewhere. Their composition $\Phi_d$ is the climax of Vol III
and is CONJECTURAL at functor level; PROVED at value level for
$d\in\{1,2\}$ and K3-fibered Class A at $d=3$.

This paper is an elementary-step derivation of Stage 1: every line is
read side-by-side with CGI--CGII, Costello--Li 2016, Gwilliam--Williams
2021, Tamarkin 2003, Kontsevich 1999. No cross-volume propagation and
no numerical claim is made beyond what each cited paper proves.

\paragraph{Conventions.} $X$ denotes a CY$_d$ manifold with
nonvanishing holomorphic volume $\Omega_X \in H^{d,0}(X)$. For
local computations we take $X = \CC^d$ with $\Omega_0 = dz_1\wedge\cdots
\wedge dz_d$. $\fg$ is a finite-dimensional Lie algebra with
invariant non-degenerate Killing pairing $\kappa_\fg(-,-)$; we always
write $\langle -, - \rangle$ for $\kappa_\fg$ extended to $\fg$-valued
Dolbeault forms. Cohomological degree $|{-}|$ is ghost degree plus
Dolbeault form degree, shifted by $-1$. $\hbarb$ is the BV quantisation
parameter. Multiplication in the graded Lie bracket on $\Omega^{0,\bullet}(X,\fg)$
uses the Koszul sign $(-1)^{|a||b|}$ on swap. The Bochner--Martinelli
kernel on $\CC^d$ is
\[
  P_{\BM}(z,w)
  \;=\; \frac{(d-1)!}{(2\pi i)^d}
        \sum_{k=1}^{d} (-1)^{k-1}
        \frac{\overline{z_k - w_k}}{\|z-w\|^{2d}}
        \,\widehat{d\bar z_k} \wedge dw_1\wedge\cdots\wedge dw_d,
\]
where $\widehat{d\bar z_k}$ means the product of $d\bar z_j$ for
$j\neq k$. At $d=3$ this is the fundamental solution of $\bar\partial$
from Bochner--Martinelli--Koppelman (Krantz 2001, \emph{Function Theory
of Several Complex Variables}, Thm~4.1.5).

\section*{Cycle 1. Attack --- The classical BV datum is not the factorisation algebra}

\begin{attack}
Six attacks on the assertion ``classical hCS has an
$\mathbb{E}_d$-holomorphic FA structure.''

\emph{A1.1 (Classical-quantum conflation).} The classical action
\(
S_{\mathrm{cl}}(\cA)=\int_X \Omega_X\wedge
\langle \cA, \bar\partial\cA + \tfrac{1}{3}[\cA,\cA]\rangle
\)
is a functional. A functional is not a factorisation algebra. The FA
structure lives on the quantum observables
$\Obs^{\mathrm{q}}(X)$, not on the action. To claim the classical data
produces Stage 1 without invoking BV quantisation is category error.

\emph{A1.2 (Shift law sign).} The shift law
$\mathrm{shift} = d-4$ produces $(d, \mathrm{shift}, \mathbb{E}_n)$
$\in\{(2,-2,\mathbb{E}_2), (3,-1,\mathbb{E}_1), (4,0,\mathbb{E}_0),
(5,+1,\mathbb{E}_5^{\mathrm{Poisson}})\}$. But the FA receives
$\mathbb{E}_d$-structure, not $\mathbb{E}_n$ with $n$ equal to the
shifted-symplectic AKSZ label. The two $\mathbb{E}$-labels measure
different things. At $d=3$ the AKSZ label is $\mathbb{E}_1$
(reflecting the $(-1)$-shifted symplectic); the FA label is
$\mathbb{E}_3$ (reflecting three holomorphic directions after
$\bar\partial$-cohomology). Both are correct; conflating them is not.

\emph{A1.3 (Ghost-number degree accounting).} Assigning ghost number
to $c$ = $-1$, $A^{(0,1)}$ = 0, $A^{*(0,2)}$ = +1, $c^{*(0,3)}$ = +2
and claiming the classical action is of total degree $0$ requires the
volume form to supply precisely three units of Dolbeault form degree,
eliminating six units of internal shift. The arithmetic 3 (from
$\Omega_X\in\Omega^{3,0}$) plus 3 (from the three
$(0,\bullet)$-Dolbeault slots summed on the two fields and their
bracket) minus $6 = 0$ is forced; ambiguity by $\pm 2$ is
phenomenologically harmless but the sign of the Serre pairing matters.

\emph{A1.4 (Gauge Lie algebra from CY category).} The claim that
$\fg := \mathrm{RHom}_\cC(G, G)[1]$ is a Lie algebra for any generator
$G$ of $\cC$ is a category-theoretic statement requiring: (a) $G$ is a
spherical object if $\cC$ is a Fukaya category; (b) $\mathrm{RHom}_\cC$
has a cyclic $A_\infty$-structure; (c) the shift $[1]$ is the standard
Lie shift. Generic CY categories admit no canonical generator; Orlov's
split-generation theorem (Orlov 2009, \emph{Trends Math}) provides one
up to Morita equivalence, but the associated Lie algebra depends on
Morita class. This is \emph{not} a canonical choice.

\emph{A1.5 (Compactness versus non-compactness of X).} Costello 2011
Thm~13.4.1 assumes compact support of the classical interaction; for
non-compact $X = \CC^d$ (where our flat-space avatar lives) one must
invoke compactly-supported Costello--Gwilliam (CGII~Prop~8.2.1) or
Costello--Li 2016 (Prop~5.2, propagator with compact support). Writing
$\Obs^{\mathrm{q}}(\CC^d)$ without specifying the compactness regime
is ambiguous.

\emph{A1.6 (Serre duality sign convention).} The pairing
$\langle A^{(0,1)}, A^{*(0,2)}\rangle$ under $\Omega_X\wedge$ gives a
$(3,3)$-form. Serre duality
$H^{0,q}(X,\fg)^\vee \simeq H^{d,d-q}(X,\fg^\vee)$ produces the
$(-1)$-shifted symplectic on $\Omega^{0,\bullet}(X,\fg)[1]$ only at
$d=3$ (shift $d-4=-1$); at $d\neq 3$ the pairing is differently
shifted, and the FA label differs accordingly. The \emph{sign} of
$\omega_{\BV}$ is fixed by the orientation of $\Omega_X$; reversing the
orientation reverses the sign of the Poisson bracket, which changes
the sign of the $\mathbb{E}_3$-algebra structure.
\end{attack}

\begin{heal}[Cycle 1: BV quantisation and the factorisation algebra of observables]
We separate the six attacks.

\emph{H1.1 (BV quantisation, not classical action, is Stage 1).} We
state the construction of Stage 1 as the BV-quantised factorisation
algebra of observables of $\hCS(d)$ on $X$, coupled to the cyclic
$A_\infty$-category $\cC$ through its endomorphism Lie algebra.
Classical data is \emph{input}; Stage 1 is the \emph{quantum output}.

\begin{definition}[Stage 1 construction, $d=3$ flat-space base case]
\label{def:stage1-flat}
Fix a finite-dimensional Lie algebra $\fg$ with invariant non-degenerate
pairing $\kappa_\fg$. The classical BV theory $\hCS(3)$ on $(\CC^3, \Omega_0)$
has:
\begin{itemize}[leftmargin=2em]
  \item Field sheaf $\cE = \Omega^{0,\bullet}(\CC^3,\fg)[1]$, as a
        sheaf of graded vector spaces on $\CC^3$.
  \item Classical BRST differential $Q_{\mathrm{cl}} = \bar\partial +
        \{\tfrac{2}{3}\langle\cA,[\cA,\cA]\rangle, \cdot\}$, i.e.\
        $Q_{\mathrm{cl}}\cA = \bar\partial\cA + \tfrac{1}{2}[\cA,\cA]$.
  \item $(-1)$-shifted symplectic pairing
        $\omega_{\BV}(\cA_1, \cA_2) = \int_{\CC^3} \Omega_0 \wedge
        \langle \cA_1, \cA_2\rangle$ on compactly supported
        $\cA_i \in \cE_c(\CC^3)$.
\end{itemize}
Apply Costello's BV machinery (\emph{Renormalization and Effective
Field Theory}, Chapters~2--5, 13): choose a gauge-fixing operator
$\bar\partial^*$ (formal adjoint of $\bar\partial$ with respect to a
flat hermitian metric on $\CC^3$), a scale-family of propagators
$P_L^\varepsilon$ regularising $\bar\partial^{-1}$, and an effective
interaction $I[L]$ solving the quantum master equation
$Q_{\mathrm{cl}}I[L] + \hbarb \Delta_L I[L] + \tfrac{1}{2}\{I[L], I[L]\}_L = 0$
at every scale $L > 0$. The quantum observables are the completed
symmetric algebra on the field dual
\[
  \Obs^{\mathrm{q}}(\cU)
  \;=\; \bigl(\Sym(\cE(\cU)^\vee[1])[[\hbarb]], \; Q_{\mathrm{cl}} + \hbarb\Delta\bigr),
  \qquad \cU\subset\CC^3 \text{ open},
\]
with $\Delta$ the BV Laplacian constructed from $P_L^\varepsilon$. The
associated prefactorisation cosheaf $\cU\mapsto\Obs^{\mathrm{q}}(\cU)$
is the Stage-1 output; we write $\Phi^{\FA}_3(\fg, \Omega_0) =
\Obs^{\mathrm{q}}$.
\end{definition}

\emph{H1.2 (Shift law vs FA label are independent).} We record the
two labels as distinct invariants.

\begin{proposition}[AKSZ shift law and FA label]
\label{prop:shift-fa}
Let $X$ be CY$_d$. The AKSZ shift of the BV pairing on $\hCS(d)$ is
$d - 4$:
\[
  \omega_{\BV}\colon \Sym^2 \cE \to \CC[d-4],
\]
which at $d = 3$ is $(-1)$-shifted symplectic, at $d = 4$ is
$0$-shifted (ordinary Poisson), at $d = 5$ is $(+1)$-shifted
(Calaque--Pantev--To\"en--Vaqui\'e--Vezzosi~2017, hereafter CPTVV).
Independently, the factorisation algebra of observables on the
underlying $(2d)$-real-dimensional manifold carries an
$\mathbb{E}_{2d}$-algebra structure \emph{before} holomorphic reduction
(CGII Theorem~5.0.1); after passing to $\bar\partial$-cohomology on
the Dolbeault complex $\Omega^{0,\bullet}$, the structure reduces to
$\mathbb{E}_d$ (Williams~2017, \emph{The Holomorphic $\sigma$-Model},
Ch.~4). The two labels are:
\[
  (\text{AKSZ shift}, \; \mathbb{E}\text{-label after Dolbeault reduction})
  \;=\; (d - 4, \; \mathbb{E}_d), \qquad d\in\{2,3,4,5\}.
\]
\end{proposition}

\emph{H1.3 (Ghost-degree computation, explicit).} On the $(0,q)$-component
of $\cA$ of ghost-shifted degree $q - 1$, the classical action
$S_{\mathrm{cl}} = \int_X \Omega_X \wedge \langle \cA, \bar\partial\cA
+ \tfrac{1}{3}[\cA,\cA]\rangle$ requires total cohomological degree
zero plus total Dolbeault form degree $d = 3$. Distributing
$\sum_{q_1, q_2} \langle\cA^{(0,q_1)}, \bar\partial \cA^{(0,q_2)}\rangle$
with $q_2 = q_1 + 1$: the non-vanishing pairs are
$(q_1, q_2) \in \{(0,1), (1,2), (2,3)\}$, contributing Dolbeault
bidegrees $(0, 1) + (0, 2) + (0, 3)$. Combined with $\Omega_X$ of
Dolbeault bidegree $(3, 0)$, total form bidegree is $(3, 3)$, a top
form; cohomological degree is $(q_1 - 1) + (q_2 - 1) + 1 = q_1 + q_2 - 1
\in\{0, 2, 4\}$, shifted by $\Omega_X\wedge$ (weight $0$ in cohomological
degree) and integrated against the BV pairing (weight $-1$). Net: the
$(0,0)+(0,1)$ pair gives $0$; $(0,1)+(0,2)$ gives $2$; $(0,2)+(0,3)$
gives $4$. Only the middle term $(0,1)+(0,2)$ is the kinetic term at
cohomological degree~0 for the gauge field; the others are the
ghost--anti-field pairings at degree~0 in BV-graded sense after the
$(-1)$-shift. \emph{No arithmetic mismatch.}

\emph{H1.4 (Canonical Lie algebra extraction: via derived
endomorphisms).} We do not require $\cC$ to be $\fg$-gauge for a
specified $\fg$; instead, Stage 1 sends $\cC$ to the ``matter-coupled''
$\hCS(d)$ whose field sheaf is $\cE_{\cC} := \Omega^{0,\bullet}(X)\otimes
\HH^\bullet(\cC,\cC)[d]$. The shift $[d]$ is the CY-shift. The
Hochschild complex carries a Gerstenhaber algebra structure
(Gerstenhaber--Voronov 1995, \emph{J.~Pure Appl.~Algebra} 103); for a
CY$_d$-category, KT-formality (Attack 3 below) upgrades this to an
$\mathbb{E}_d$-algebra. Under factorisation-homology reduction on
$X = \CC^d$, one recovers an $\mathbb{E}_d$-algebra structure on the
observables; when $\cC = \Perf(\text{pt})$ (equivalently $\fg = 0$),
the observables collapse to the structure sheaf.

\emph{H1.5 (Compactness regime fixed by choice of propagator
support).} We fix the regime: $\Phi^{\FA}_d(\cC)$ is constructed using
compactly-supported propagators (CGII Ch.~8--10) on arbitrary
(compact or non-compact) CY$_d$ $X$. For compact $X$, Costello~2011
Thm~13.4.1 applies directly; for non-compact $X$, CGII~Prop~8.2.1
(locality of compactly-supported observables) replaces Thm~13.4.1.

\emph{H1.6 (Orientation of $\Omega_X$ fixed once and for all).} We
fix $\Omega_X = dz_1\wedge dz_2\wedge dz_3$ on the flat-space base
case and transport to curved $X$ by the complex structure. The sign
of $\omega_{\BV}$ is then fixed.
\end{heal}

\section*{Cycle 2. Attack --- The $\mathbb{E}_d$-structure is not automatic}

\begin{attack}
Four attacks on the assertion ``$\Obs^{\mathrm{q}}(\CC^d)$ has an
$\mathbb{E}_d$-structure.''

\emph{A2.1 (Which $\mathbb{E}$?).} There are four distinct
$\mathbb{E}$-structures floating near the BV FA on $\CC^d$:
\begin{itemize}[leftmargin=2em]
  \item[(i)] Topological $\mathbb{E}_{2d}$: from $(2d)$-real-dimensional
             translation invariance on $\RR^{2d}$, before any
             holomorphic consideration.
  \item[(ii)] Holomorphic $\mathbb{E}_d$: after choosing a complex
             structure, restricting to holomorphic translation
             invariance, and taking $\bar\partial$-cohomology on the
             Dolbeault complex. This is Williams 2017, Costello--Li 2016.
  \item[(iii)] Chiral $\mathbb{E}_1$ on a curve: after further
              specialising the holomorphic FA to a curve $C \subset X$
              via factorisation homology over a transverse $(d-1)$-cycle.
              This is Stage 2 and is not our concern here.
  \item[(iv)] Derived $\mathbb{E}_{?}$ on the Drinfeld double of $\cC$:
              a distinct structure appearing in the $\cC\to G(\cC) =
              D(Y^+(\cC))$ assignment; \emph{not} what Stage 1
              constructs.
\end{itemize}
Calling ``the'' $\mathbb{E}_d$-structure without specifying which is
ambiguous. (This is the Vol I cache conflation catalogued at
\texttt{feedback\_e1\_chiral\_multiple\_notions.md}.)

\emph{A2.2 (Dunn additivity: why should it hold?).} The claim
\[
  \mathbb{E}_d = \mathbb{E}_1 \otimes_D \mathbb{E}_1 \otimes_D \cdots \otimes_D \mathbb{E}_1
  \qquad (d \text{ factors})
\]
is Dunn--Lurie additivity (Lurie, \emph{Higher Algebra} 5.1.2.2). It
applies to locally-constant translation-equivariant factorisation
algebras on $\RR^d$. On $\CC^d$ (viewed as $\RR^{2d}$), naively Dunn
additivity would give $\mathbb{E}_{2d}$ (one $\mathbb{E}_1$ per real
direction), not $\mathbb{E}_d$. The reduction from $\mathbb{E}_{2d}$
to $\mathbb{E}_d$ requires a non-trivial Mori-invariant statement:
that the FA is ``$\bar\partial$-rigid'' in the sense of Williams 2017
Thm~4.4.1 (each holomorphic direction $z_k$ contributes a single
$\mathbb{E}_1$-factor post-Dolbeault cohomology, not two). Cache entry
450 / AP932 in Vol I and AP-V2-36 in Vol II flag this as a common
conflation.

\emph{A2.3 (The Dolbeault cohomology step is model-dependent).}
Passing from the $\mathbb{E}_{2d}$-structure on $\Obs^{\mathrm{q}}$ to
the $\mathbb{E}_d$-structure on $H^\bullet_{\bar\partial}\Obs^{\mathrm{q}}$
requires a \emph{model} for the holomorphic FA: the Costello--Li
Dolbeault-cohomology model makes the operation literal, but other
models (e.g.\ \v{C}ech--Dolbeault, or Koszul dualisation) need the
comparison theorems that these models are all quasi-isomorphic
(Dupont 1976; Costello--Gwilliam Vol II, Ch.~5).

\emph{A2.4 (Williams 2017 is one among several proofs).} Independent
proofs of the $\mathbb{E}_d$-structure on
$H^\bullet_{\bar\partial}\Obs^{\mathrm{q}}(\CC^d)$ include:
Gwilliam--Williams 2021 (direct construction via Bochner--Martinelli
propagator; Thm~1.1); Costello--Li 2016 (renormalisation-based argument;
Prop~6.3); Williams 2020 \emph{J.~Pure Appl.~Algebra} 224
(factorisation-homology argument). Claiming ``the'' proof without
naming which is imprecise.
\end{attack}

\begin{heal}[Cycle 2: the four $\mathbb{E}$-structures, isolated and labelled]

\emph{H2.1 (Labelled four $\mathbb{E}$-structures).} We fix the
notation:
\begin{itemize}[leftmargin=2em]
  \item $\Obs^{\mathrm{top}}$: topological FA on $\RR^{2d}$, carrying
        $\mathbb{E}_{2d}$ by CGII~Thm~5.0.1.
  \item $\Obs^{\mathrm{hol}}$: holomorphic FA on $\CC^d$, which is
        $\Obs^{\mathrm{top}}$ \emph{after imposing holomorphic
        translation invariance}; carries an
        $\mathbb{E}_{2d}$-action whose restriction to the antiholomorphic
        generators is invertible (i.e., the FA is locally constant in
        antiholomorphic directions).
  \item $\Obs^{\FA}_d := H^\bullet_{\bar\partial} \Obs^{\mathrm{hol}}$:
        Dolbeault-cohomology FA on $\CC^d$, carrying $\mathbb{E}_d$ by
        Dunn additivity applied to the reduced model.
  \item $\Obs^{\mathrm{chir}}_1 := \mathrm{Sp}^{\mathrm{ch}}_{\Sigma_{d-1}, C}
        \Obs^{\FA}_d$: Stage 2 output, an $\mathbb{E}_1$-chiral algebra on $C$.
\end{itemize}
Stage 1 produces \emph{$\Obs^{\FA}_d$}. The $\mathbb{E}_d$-label is
unambiguous once we fix the Dolbeault-reduction step.

\emph{H2.2 (Dunn additivity on the Dolbeault complex).}

\begin{proposition}[Holomorphic Dunn additivity, Williams 2017
Thm~4.4.1]
\label{prop:hol-dunn}
Let $F$ be a locally-constant translation-equivariant factorisation
algebra on $\CC^d$, equipped with an action of the antiholomorphic
translation group $\CC^d$ that is locally constant (each
antiholomorphic translation is a quasi-isomorphism). Then
$H^\bullet_{\bar\partial} F$ carries a locally-constant translation-equivariant
$\mathbb{E}_d$-factorisation algebra structure, such that the natural
inclusion $H^\bullet_{\bar\partial} F \hookrightarrow F$ is a quasi-iso
of $\mathbb{E}_d$-algebras (with $F$ viewed as $\mathbb{E}_d$ by
restricting $\mathbb{E}_{2d}$ along the holomorphic embedding
$\mathbb{E}_d \hookrightarrow \mathbb{E}_{2d}$). By Dunn additivity
(Lurie HA 5.1.2.2) on the Dolbeault-reduced model,
\[
  \Obs^{\FA}_d \;\simeq\; \underbrace{\Obs^{\FA}_1 \boxtimes_D \cdots \boxtimes_D \Obs^{\FA}_1}_{d \text{ factors}}
\]
where $\Obs^{\FA}_1$ is the $\mathbb{E}_1$-holomorphic FA on $\CC$
(two-dimensional chiral CFT observables).
\end{proposition}

\begin{proof}[Sketch]
Each holomorphic direction $z_k$ contributes a 1-parameter family of
open-disc embeddings; the Dolbeault-reduction step identifies the
antiholomorphic translations as identity maps modulo exact, so each
holomorphic coordinate contributes an $\mathbb{E}_1$-factor. The
product over $k = 1, \ldots, d$ is Dunn tensor product. See
Williams 2017 \S4.4; Gwilliam--Williams 2021 \S3 for the case $d=3$.
\end{proof}

\emph{H2.3 (Model comparison).} The Dolbeault-cohomology model, the
\v{C}ech--Dolbeault model, and the Koszul-dual model are all
quasi-isomorphic as $\mathbb{E}_d$-algebras (Costello--Gwilliam
Vol II, Ch.~5 Prop~5.2.1). We use the Dolbeault-cohomology model for
definiteness.

\emph{H2.4 (Three proof paths identified).}

\begin{proposition}[Three independent proofs of Stage-1 $\mathbb{E}_d$-structure]
\label{prop:three-proofs}
Let $X = \CC^d$ with flat CY structure. Three proofs of the
$\mathbb{E}_d$-structure on $\Phi^{\FA}_d(\cC)$ are available:
\begin{itemize}[leftmargin=2em]
  \item[(P1)] \emph{Gwilliam--Williams 2021 direct.} Construct the
              Bochner--Martinelli propagator on $\CC^d$ explicitly; the
              prefactorisation cosheaf $\Obs^{\mathrm{q}}$ constructed
              from this propagator, plus the associated Dolbeault
              reduction, is proved to carry an $\mathbb{E}_d$-structure
              in Thm~1.1 loc.~cit.
  \item[(P2)] \emph{Costello--Li 2016 renormalisation.} Apply the
              renormalisation-group machinery of Costello 2011 with
              the Bochner--Martinelli gauge-fixing; the effective BV
              quantisation obeys the hypothesis of CGII Thm~5.0.1
              (holomorphic translation equivariance) after Dolbeault
              reduction, giving $\mathbb{E}_d$-structure.
  \item[(P3)] \emph{Factorisation-homology argument, Williams 2020.}
              Identify $\Phi^{\FA}_d(\cC)$ with the factorisation
              homology $\int_{\CC^d} \cA$ of an associated
              $\mathbb{E}_d$-algebra $\cA$ (constructed from $\cC$ via
              KT-formality); by the factorisation-homology axioms, the
              result is an $\mathbb{E}_d$-FA.
\end{itemize}
All three paths produce the same FA up to contractible choice.
\end{proposition}
\end{heal}

\section*{Cycle 3. Attack --- Kontsevich--Tamarkin formality: who proves what, for which $d$?}

\begin{attack}
Five attacks on the invocation ``KT formality.''

\emph{A3.1 (Which $d$ exactly?).} Kontsevich 1999 (\emph{Deformation
quantization of Poisson manifolds}, \emph{Lett.~Math.~Phys.} 66)
proves $\mathbb{E}_2$-formality for $C^\infty(M)$: the Hochschild
complex of functions on a smooth manifold, equipped with its natural
Gerstenhaber = $\mathbb{E}_2$-algebra structure, is formal. This is
$d = 2$. Tamarkin 2003 (\emph{Action of the Grothendieck--Teichm\"uller
group on the Hochschild complex}, preprint) proves
$\mathbb{E}_2$-formality for the Hochschild complex of any
$A_\infty$-algebra, via the action of a Drinfeld associator; this is
also $d = 2$ \emph{as an operation on the Gerstenhaber structure},
not $d \geq 3$. For $d \geq 3$, the corresponding statement is the
``higher Deligne conjecture,'' proved by Tamarkin in the same paper
(Thm~2.1) and reproved by others: the little $d$-discs operad acts on
the $d$-fold iterated Hochschild complex.

\emph{A3.2 (Formality depends on an associator; is it rational or
transcendental?).} Kontsevich's proof uses the Kontsevich graph
integrals on the hyperbolic upper half-plane; the associated constants
are $\QQ$-linear combinations of multiple zeta values $\zeta(m_1, \ldots,
m_k)$. Tamarkin's proof uses a Drinfeld associator $\Phi_{\mathrm{Dr}}$
with coefficients in the $\QQ$-algebra of motivic MZVs. Both
constants live in a $\mathrm{GRT}_1(\QQ)$-torsor
(Willwacher 2014, \emph{Invent.~Math.} 200 Thm~1.2): the space of
associators is a principal homogeneous space under the prounipotent
group $\mathrm{GRT}_1(\QQ)$. Hence ``KT formality'' does not specify a
unique formality quasi-isomorphism; it specifies a torsor, and a
specific quasi-iso requires a specific associator.

\emph{A3.3 (CY$_d$ condition on $\cC$ is not innocent).} Tamarkin's
theorem requires the $A_\infty$-algebra (or category) to be
\emph{cyclic} to install the Gerstenhaber algebra structure on
Hochschild cochains. The CY$_d$-condition
$\eta\colon \HH^\bullet(\cC,\cC)[d] \xrightarrow{\sim} \HH_\bullet(\cC,\cC)^\vee$
gives this cyclicity, plus a shift $[d]$ that is essential for the
$\mathbb{E}_d$-labeling. Without CY$_d$-structure, the formality
statement is only $\mathbb{E}_2$, not $\mathbb{E}_d$.

\emph{A3.4 (Willwacher sharpening).} Willwacher 2014 (\emph{Invent.~Math.}
200 Thm~1.2) proves that the graph complex GC$^0$ has cohomology
$\cong \mathfrak{grt}_1$, so the automorphism group of the
$\mathbb{E}_2$-operad acting on Hochschild cochains is
$\mathrm{GRT}_1(\QQ)$ itself. This refines KT formality by identifying
the torsor of formality quasi-isos explicitly. For higher $d$,
Willwacher + Fresse 2020 (\emph{Compositio Math.} 156) extend this:
the automorphisms of the little $d$-discs operad act as
$\mathrm{GT}_d(\QQ)$-torsor on formality quasi-isos, with
$\mathrm{GT}_d(\QQ) \simeq \mathrm{GRT}_1(\QQ)$ for all $d \geq 2$
(Fresse--Willwacher 2020 \emph{Comp.~Math.} 156 Thm~A).

\emph{A3.5 (A CY$_d$-category is not a CY $d$-fold).} The KT-formality
assertion is about the Hochschild complex of the CY$_d$-\emph{category}
$\cC$, not about Poisson cohomology of a CY$_d$-variety $X$. The two
are related by $\cC = D^b\Coh(X)$ at $d = 3$, in which case
Hochschild cochains of $D^b\Coh(X)$ compute the polyvector-field
cohomology $H^\bullet(X, \wedge^\bullet TX)$ (Hochschild--Kostant--Rosenberg
theorem 1962; Swan 1996). In this case KT formality on the category
side is equivalent to Kontsevich formality on the variety side, but
only up to the HKR quasi-iso; conflating them without naming HKR is
an error flagged at AP-CY142 in Vol III.
\end{attack}

\begin{heal}[Cycle 3: Kontsevich--Tamarkin $\mathbb{E}_d$-formality, stated correctly]

\emph{H3.1--H3.2 (KT formality, unified statement).}

\begin{theorem}[KT formality with $\mathrm{GRT}_1(\QQ)$-torsor scope,
Kontsevich--Tamarkin 1999--2003 with Willwacher--Fresse 2014--2020
refinements]
\label{thm:kt-formality}
Let $\cC$ be a cyclic $A_\infty$-algebra over $\CC$ of
Calabi--Yau-$d$-category type, i.e.\ equipped with a non-degenerate
Mukai pairing
$\eta\colon \HH^\bullet(\cC,\cC)[d] \xrightarrow{\sim} \HH_\bullet(\cC,\cC)^\vee$
of cohomological degree $d$. The Hochschild cochain complex
$\HH^\bullet(\cC,\cC)$ carries:
\begin{itemize}[leftmargin=2em]
  \item its natural Gerstenhaber $\mathbb{E}_2$-algebra structure
        (Gerstenhaber 1963; Voronov 2000, \emph{Ann.~Sci.~ENS}),
        upgraded to an $\mathbb{E}_d$-algebra structure using the
        CY$_d$-pairing to produce a trace (Kontsevich 1999 for
        $d = 1$; Costello 2007 \emph{Adv.~Math.} 210 Thm~A for $d = 2$;
        Tamarkin 2003 Thm~2.1 for general $d$);
  \item an $\mathbb{E}_d$-formality: there is a zig-zag
        of $\mathbb{E}_d$-algebra quasi-isomorphisms
        \[
          \HH^\bullet(\cC,\cC)\otimes_\CC (\Lie[d-1]\text{-algebra})
          \;\simeq\; H^\bullet(\HH^\bullet(\cC,\cC)) \otimes_\CC (\Lie[d-1]\text{-algebra})
        \]
        depending on a choice of associator $\Phi \in \mathrm{Ass}_d(\QQ)$,
        where $\mathrm{Ass}_d(\QQ)$ is the $\mathrm{GRT}_1(\QQ)$-torsor of
        Drinfeld associators (Fresse--Willwacher 2020 Thm~A).
\end{itemize}
Transcendental constants: Kontsevich graph integrals produce a specific
element in the motivic-MZV-subring of periods; all other associators
differ by a $\mathrm{GRT}_1(\QQ)$-gauge.
\end{theorem}

\emph{H3.3 (CY$_d$-condition absorbed into the hypothesis).} The
cyclicity datum is fed in as part of the input $\cC$; without it the
theorem only gives $\mathbb{E}_2$-formality.

\emph{H3.4 (Willwacher/Fresse--Willwacher refinement).} The torsor
statement is the sharpening: the formality quasi-iso is not
canonical, but its underlying torsor structure $\mathrm{Ass}_d(\QQ)$
is canonical. Stage 1 is therefore canonical up to contractible
choice (each choice of associator defines a quasi-iso, and the space
of associators is contractible to any chosen basepoint in
$\mathrm{GRT}_1(\QQ)$-gauge).

\emph{H3.5 (HKR comparison with CY variety side, for $\cC = D^b\Coh(X)$).}

\begin{proposition}[HKR bridge at $\cC = D^b\Coh(X)$, $X$ smooth
projective CY$_3$]
\label{prop:hkr-bridge}
Let $X$ be a smooth projective CY$_3$ and $\cC = D^b\Coh(X)$. Then
\[
  \HH^\bullet(\cC,\cC) \;\simeq\; \bigoplus_{p+q = \bullet} H^p(X, \wedge^q TX)
\]
as $\mathbb{E}_3$-algebras, where the right-hand side carries the
Schouten bracket and the Cap product with $\Omega_X$ identifies
$\wedge^q TX \simeq \Omega_X^{3-q}$ as a $(-1)$-shifted symplectic
structure. This is the HKR + Kontsevich--Swan theorem (HKR 1962;
Swan 1996; Caldararu 2005, \emph{Advances in Math.} 194). Under this
identification, KT formality on the category side reduces to
Kontsevich formality on the polyvector-field side.
\end{proposition}
\end{heal}

\section*{Cycle 4. Attack --- Costello--Gwilliam--Li locality: what does it buy?}

\begin{attack}
Six attacks on the invocation ``CGL locality.''

\emph{A4.1 (Locality versus $\mathbb{E}_n$-structure: which is the theorem?).}
The statement ``the FA of observables on $X$ is determined by its
restriction to small open disks'' has two distinct mathematical
meanings:
\begin{itemize}[leftmargin=2em]
  \item[(L1)] \emph{Cosheaf locality}: $\Obs^{\mathrm{q}}$ is a
              \v{C}ech-cosheaf for the topology of disjoint open
              embeddings, i.e.\ the FA condition itself. This is
              tautological: any FA is defined by locality. CGII~Prop~6.1.3.
  \item[(L2)] \emph{Locally-constant structure}: $\Obs^{\mathrm{q}}$ is
              locally constant (each inclusion of disks in disks is a
              quasi-iso). This is the hypothesis of Lurie HA 5.3.4.10
              and gives the identification with an
              $\mathbb{E}_n$-algebra.
  \item[(L3)] \emph{Holomorphic locally-constant structure}: the
              antiholomorphic translations are quasi-isos. This is
              weaker than (L2) and gives the identification with an
              $\mathbb{E}_d$-algebra on $\CC^d$ after Dolbeault
              reduction (Williams 2017).
\end{itemize}
Stage 1 needs (L3), not the weaker (L1). The ``locality'' label is
ambiguous.

\emph{A4.2 (Costello--Li 2016: which theorem exactly?).}
Costello--Li 2016 (\emph{arXiv:1605.09930, Twisted supergravity and
its quantization}) is cited often but the specific result invoked
for Stage 1 is Proposition~5.2, which gives the existence of a
Bochner--Martinelli-style propagator on CY$_3$ with the requisite
compact support and symmetry. This is not an $\mathbb{E}_d$-statement;
it is a gauge-fixing statement. The $\mathbb{E}_d$-structure comes
from Williams 2017 or Gwilliam--Williams 2021.

\emph{A4.3 (Compactly supported versus distributional observables).}
CGII distinguishes $\Obs^{\mathrm{q}}$ (distributional) from
$\Obs^{\mathrm{q}}_c$ (compactly supported). Locality works in both
lanes but with different target categories: $\Obs^{\mathrm{q}}$ is
valued in \emph{completed} symmetric algebra (distribution space),
$\Obs^{\mathrm{q}}_c$ in the compactly supported symmetric algebra
(Schwartz-space). The $\mathbb{E}_d$-structure exists in both but on
different ambient categories.

\emph{A4.4 (``Canonical up to contractible choice'' — of what?).}
The phrase means: the space of BV renormalisations (gauge fixings,
propagators, counterterm schemes) is contractible; any two choices
produce equivalent FAs. Costello 2011 \S6--9 proves this for the
scale-family of propagators $P_L^\varepsilon$ (any two
$L_1 < L_2$ give isomorphic effective actions via RG flow). For the
gauge-fixing, contractibility is the statement that the space of
gauge-fixing conditions $\bar\partial^*$ is an affine bundle over a
contractible base (the space of hermitian metrics on $X$). The
``contractible choice'' language requires these two contractibility
statements as separate inputs, not one.

\emph{A4.5 (Li--Li 2020, Li 2020: the $\hCS$-specific refinement).}
For the $\hCS$-avatar specifically, the existence and uniqueness of
the BV quantisation up to contractible choice require:
Li~2020 \emph{J.~Differential Geom.} 114 (Feynman-diagram
regularisation on CY$_3$); Li--Li~2020 \emph{Commun.~Math.~Phys.} 376
(renormalisation-group analysis of $\hCS$). These are more specific
than general CGII.

\emph{A4.6 (The ``canonical'' is only up to contractible choice in
the sense of $\infty$-categorical equivalence).} ``Canonical'' in
the $(\infty,1)$-sense means the space of constructions is connected
and simply-connected; the construction is not a functor valued in a
1-category but in the $(\infty,1)$-category of
$\mathbb{E}_d$-holomorphic FAs.
\end{attack}

\begin{heal}[Cycle 4: CGL locality, stated tri-layered and crisply]

\emph{H4.1 (Three locality statements, all invoked).}

\begin{theorem}[CGL locality, tri-layered, CGII + Costello--Li 2016]
\label{thm:cgl-locality}
Let $X$ be a complex manifold of complex dimension $d$ and let
$\Phi^{\FA}_d(\cC)$ be the quantum observables of $\hCS(d)$ on $X$
coupled to CY$_d$-category $\cC$, constructed as in
Definition~\ref{def:stage1-flat}. Then:
\begin{itemize}[leftmargin=2em]
  \item[(L1)] \emph{Cosheaf locality.} $\Phi^{\FA}_d(\cC)$ is a
              prefactorisation cosheaf: for every finite family of
              pairwise disjoint open sets
              $\cU_1,\ldots,\cU_k\subset\cU\subset X$ and the
              corresponding structure map
              $\Phi^{\FA}_d(\cC)(\cU_1) \otimes \cdots \otimes \Phi^{\FA}_d(\cC)(\cU_k)
              \to \Phi^{\FA}_d(\cC)(\cU)$, the FA condition holds
              (CGII~Prop~6.1.3).
  \item[(L2)] \emph{Locally-constant-ness in holomorphic direction.}
              For any open inclusion $\cU\hookrightarrow\cU'$ that is a
              holomorphic tubular retract (i.e.\ $\cU = \cU'\setminus
              (\text{antiholomorphic bundle})$), the structure map
              $\Phi^{\FA}_d(\cC)(\cU) \to \Phi^{\FA}_d(\cC)(\cU')$ is a
              quasi-isomorphism. This is the content of
              Costello--Li~2016 Prop~5.2 (Bochner--Martinelli
              propagator kills the antiholomorphic direction) combined
              with CGII~Thm~5.0.1 (translation equivariance of the
              BV-quantised observables).
  \item[(L3)] \emph{Descent to Dolbeault-cohomology FA.} The
              Dolbeault-reduced FA
              $\Phi^{\FA}_d(\cC)^{\bar\partial} := H^\bullet_{\bar\partial}
              \Phi^{\FA}_d(\cC)$ is a locally-constant
              translation-equivariant factorisation algebra on
              $\CC^d$, hence by Lurie HA 5.3.4.10 equivalent to an
              $\mathbb{E}_d$-algebra, the $\mathbb{E}_d$-holomorphic FA
              label of Stage 1. (Williams 2017 Thm~4.4.1; Gwilliam--Williams
              2021 Thm~1.1.)
\end{itemize}
Canonicity: the construction is canonical up to contractible choice
of gauge-fixing operator $\bar\partial^*$ (the space of
compact-support compatible gauge-fixing operators is contractible by
Costello 2011 \S6.5) and scale-family propagator $P_L^\varepsilon$
(the RG-flow space is contractible by Costello 2011 Thm~8.6.1).
\end{theorem}

\emph{H4.2 (Costello--Li 2016 Prop~5.2, explicit reference).}

\begin{proposition}[Bochner--Martinelli propagator, Costello--Li 2016
Prop~5.2]
\label{prop:bm-propagator}
On $\CC^3 = \{(z_1, z_2, z_3)\}$ equipped with the flat hermitian
metric, the Bochner--Martinelli kernel $P_{\BM}(z, w)$ of
\S\emph{Conventions} is a distributional bisolenoidal $(1,2)$-form on
$\CC^3 \times \CC^3$ minus the diagonal, satisfying
\[
  \bar\partial_z P_{\BM}(z, w) = \delta_\Delta(z - w) \cdot \Omega_0(w),
  \qquad
  \bar\partial_w P_{\BM}(z, w) = \delta_\Delta(z - w) \cdot \Omega_0(z).
\]
Equivalently, $P_{\BM}$ is the kernel of $\bar\partial^{-1}$ on the
orthogonal complement of harmonic $(0,\bullet)$-forms with compact
support, modulo the hermitian metric gauge. The regularisation
$P_L^\varepsilon$ at scales $\varepsilon < L$ is the smoothed version
obtained by cutting off the diagonal singularity with a smooth
bump-function family; contractibility of the space of such families
is Costello 2011 Prop~6.5.3.
\end{proposition}

\emph{H4.3 (Compact vs distributional lanes fixed).} We work in the
distributional lane: $\Obs^{\mathrm{q}}$ is valued in completed
symmetric algebra of distributions, following CGII~\S3.1 and Williams
2017. The compactly supported lane $\Obs^{\mathrm{q}}_c$ appears in
Stage 2 via factorisation homology pairings.

\emph{H4.4--H4.6 (Contractible-choice canonicity, unpacked).}

\begin{proposition}[Canonicity of Stage 1 up to contractible choice]
\label{prop:canonicity-stage1}
The Stage-1 FA $\Phi^{\FA}_d(\cC)$ is a functor
\[
  \Phi^{\FA}_d\colon \mathrm{CY}_d\text{-Cat}^{\mathrm{cyc},\mathrm{prop}}
  \longrightarrow \mathbb{E}_d\text{-}\HolFA(X)
\]
of $(\infty,1)$-categories, canonical up to contractible choice of:
\begin{itemize}[leftmargin=2em]
  \item the associator $\Phi \in \mathrm{Ass}_d(\QQ)$ controlling
        KT-formality (space of associators is a $\mathrm{GRT}_1(\QQ)$-torsor,
        contractible in $\QQ$-algebra of periods; Fresse--Willwacher 2020);
  \item the gauge-fixing operator $\bar\partial^*$ (space of
        compactly supported gauge-fixing operators is contractible;
        Costello 2011 Prop~6.5.3);
  \item the scale-family of propagators $P_L^\varepsilon$
        (RG-flow space is contractible; Costello 2011 Thm~8.6.1).
\end{itemize}
All three spaces of choices are independently contractible; their
product is contractible. The functor $\Phi^{\FA}_d$ lifts to the
$(\infty,1)$-category of contractibly-chosen data.
\end{proposition}
\end{heal}

\section*{Cycle 5. Attack --- The $d=3$ case: 6D hCS is the canonical avatar?}

\begin{attack}
Seven attacks on the assertion ``6D holomorphic Chern--Simons on
$\CC^3$ is the canonical $\mathbb{E}_3$-holomorphic FA for $\mathrm{CY}_3$
categories.''

\emph{A5.1 (Why hCS and not some other theory?).} Many holomorphic
theories on CY$_3$ produce observables that are $\mathbb{E}_3$-FAs:
holomorphic BF, holomorphic $\sigma$-models, twisted M5-brane
worldvolume theory, twisted $(2,0)$ theory. Why is $\hCS$ ``the''
canonical one?

\emph{A5.2 (Gauge Lie algebra dependence).} $\hCS$ depends on a
choice of finite-dimensional Lie algebra $\fg$ with invariant pairing.
A general CY$_3$-category $\cC$ is infinite-dimensional (e.g.\
$D^b\Coh(X)$ for $X$ smooth projective CY$_3$). The identification
$\fg = \HH^{\bullet}(\cC,\cC)[3]$ makes $\fg$ the Hochschild complex,
which is an $A_\infty$-Lie-algebra (not an ordinary Lie algebra). The
$\hCS$ action needs a Lie-algebra structure; Koszul-dualising the
infinity-Lie-algebra to a plain one requires a (non-canonical) minimal
model.

\emph{A5.3 (Matter-coupling versus gauge-coupling).} Two distinct
couplings of $\cC$ to $\hCS$:
\begin{itemize}[leftmargin=2em]
  \item[(C1)] \emph{Gauge-coupling}: fix $\fg$ and realise $\cC$ as a
              $\fg$-module category; the $\hCS$ field space has fiber
              $\fg$.
  \item[(C2)] \emph{Matter-coupling}: introduce a matter field sheaf
              $\cM = \Omega^{0,\bullet}(X)\otimes \cC^{\mathrm{fiber}}$
              coupled to the $\hCS$ gauge field $\cA$ via a cubic
              interaction; the fiber carries $\cC$-module structure.
\end{itemize}
Stage 1 uses (C2) when $\cC$ is not already $\fg$-gauge for a specific
$\fg$. The distinction matters for the $\mathbb{E}_d$-structure (both
produce $\mathbb{E}_d$-FAs, but via different operadic routes).

\emph{A5.4 (Gwilliam--Williams 2021 computes on $\CC^3$; what about
curved $X$?).} Gwilliam--Williams 2021 explicitly computes the
$\mathbb{E}_3$-algebra of observables for $\hCS$ on $(\CC^3, \Omega_0,
\fg)$. For curved CY$_3$ $X$, the FA is constructed locally using
the flat-space model and glued via the CY$_3$-structure; the
$\mathbb{E}_3$-structure is defined globally, but the explicit
Bochner--Martinelli propagator is local.

\emph{A5.5 (Deformation theory = $\HH^\bullet_{\mathbb{E}_3}[3]$).}
Costello~2013 (\emph{Sci.~Bull.~Math.} 137) proves the deformation
complex of 6D $\hCS$ on $\CC^3$ is
$\HH^\bullet_{\mathbb{E}_3}(\Obs^{\FA}_3(\cC), \Obs^{\FA}_3(\cC))[3]$;
this identifies deformations with Hochschild cohomology of the
$\mathbb{E}_3$-algebra. For $\fg$ simple with Killing form,
$\HH^0_{\mathbb{E}_3} = \CC \cdot \mathrm{Kil}$, matching the
Yangian-deformation slot of Costello--Witten--Yamazaki 2017.

\emph{A5.6 (Costello--Dimofte--Paquette 2020: factorisation algebras
and chiral CFT).} For $\fg = \widehat{\mathfrak{gl}}_1$, the
all-orders-in-$\hbarb$ non-abelian $\hCS$ observable algebra is the
Yangian $Y_{\epsilon_1,\epsilon_2,\epsilon_3}(\widehat{\mathfrak{gl}}_1)$
(Costello--Dimofte--Paquette 2020 Thm~1.1). This is the proof that
$\hCS$ is non-trivial at the quantum level for abelian gauge algebras.

\emph{A5.7 (Relation to $\CoHA(\CC^3)$).} At $\fg = \widehat{\mathfrak{gl}}_1$,
Schiffmann--Vasserot 2013 (\emph{Duke Math.~J.} 162) constructs the
cohomological Hall algebra $\CoHA(\CC^3) = Y^+(\widehat{\mathfrak{gl}}_1)$.
Vol III identifies this with $\Phi^{\FA}_3(\cC) = \Phi^{\FA}_3(\CoHA(\CC^3))$
applied to the perverse-sheaf category on quiver moduli; the
$\mathbb{E}_3$-algebra structure on $Y^+(\widehat{\mathfrak{gl}}_1)$
\emph{is} the $\mathbb{E}_3$-FA at the $(\CC^3, \Omega_0)$ base point.
\end{attack}

\begin{heal}[Cycle 5: 6D hCS is canonical as the $\hCS(3)$ avatar of Stage 1]

\emph{H5.1 (``Canonical'' means ``unique up to the contractible-choice
$(\infty,1)$-categorical equivalence'').} We do not claim $\hCS$ is
the unique CY$_3$-theory producing $\mathbb{E}_3$-FAs; we claim Stage 1 is
canonically isomorphic to $\hCS$ observables at the $\mathrm{CY}_3$-level
because KT formality + CGL locality specify the $\mathbb{E}_3$-FA
up to contractible choice, and the $\hCS$ realisation at $\fg = $
endomorphism Lie algebra of a generator is one representative. Other
representatives (holomorphic $BF$, twisted M5, etc.) differ by
quasi-isos.

\begin{theorem}[6D hCS as canonical $d=3$ avatar, Costello~2013 +
Gwilliam--Williams~2021 + Costello--Dimofte--Paquette~2020]
\label{thm:6dhcs-canonical-avatar}
Let $\cC$ be a cyclic $A_\infty$-category of CY$_3$ type with a
compact generator $G$ and $\fg := \mathrm{RHom}_\cC(G,G)[1]$ its
endomorphism Lie algebra (well-defined up to Morita equivalence by
Orlov 2009). The Stage-1 $\mathbb{E}_3$-holomorphic factorisation
algebra
\[
  \Phi^{\FA}_3(\cC) \;=\; \Obs^{\mathrm{q}}_{\hCS(3)\text{ on }(\CC^3,\Omega_0,\fg)}
\]
is canonically isomorphic (in the $(\infty,1)$-category of
$\mathbb{E}_3$-holomorphic FAs on $\CC^3$, up to contractible choice)
to the BV-quantised observables of six-dimensional holomorphic
Chern--Simons theory on $(\CC^3, \Omega_0, \fg)$. Under this
identification:
\begin{itemize}[leftmargin=2em]
  \item For $\cC = \Coh(\mathrm{pt})$, $\fg = 0$: $\Phi^{\FA}_3 = \CC$
        (trivial).
  \item For $\cC = \Coh(E)$ (elliptic curve, CY$_1$): not directly
        Stage 1; Vol III constructs this via different route at $d = 1$.
  \item For $\cC = D^b\Coh(X)$ with $X$ smooth projective CY$_3$:
        $\fg$ is the derived polyvector fields $\bigoplus_q
        H^{\bullet}(X,\wedge^q TX)$; $\Phi^{\FA}_3(\cC)$ is the
        observable FA of $\hCS$ with this derived Lie algebra.
  \item For $\cC = \CoHA(\CC^3)$ (quiver CoHA on $\CC^3$):
        $\Phi^{\FA}_3(\cC) = Y^+(\widehat{\mathfrak{gl}}_1)$ as
        $\mathbb{E}_3$-algebra, matching Schiffmann--Vasserot 2013
        (Costello--Dimofte--Paquette 2020 Thm~1.1).
\end{itemize}
\end{theorem}

\emph{H5.2 (Gauge Lie algebra: $A_\infty$-Lie algebra directly, no
strictification).} We work in the $(\infty,1)$-categorical lane: the
``Lie algebra'' input to $\hCS$ is an $L_\infty$-algebra, and
$\HH^\bullet(\cC,\cC)[3]$ is the canonical such input. Strict Lie
algebras are a special case where the higher brackets vanish.

\emph{H5.3 (Matter-coupling recipe, explicit).}

\begin{construction}[Matter-coupled 6D hCS]
\label{cons:matter-coupled}
Given $\cC$ CY$_3$-category and $\hCS(3)$ with a chosen reference
abelian Lie algebra $\fg_0 = \CC$ (just for simplicity), introduce a
matter sheaf $\cM_\cC = \Omega^{0,\bullet}(X) \otimes_\CC \HH_\bullet(\cC)$
and a coupling interaction
\[
  S_{\mathrm{mat}}[\cA, \psi] = \int_X \Omega_X \wedge \langle \psi,
  (\bar\partial + \cA) \psi \rangle + \int_X \Omega_X \wedge \langle \psi,
  \mu_\cC(\psi) \rangle
\]
where $\mu_\cC$ is the $A_\infty$-operation from $\cC$'s module
structure and $\langle -, - \rangle$ is the Mukai pairing from the
CY$_3$-structure. The total action
$S = S_{\hCS}[\cA] + S_{\mathrm{mat}}[\cA, \psi]$ has BV quantisation
and the resulting observables are $\Phi^{\FA}_3(\cC)$.
\end{construction}

\emph{H5.4 (Curved-$X$ extension via gluing).}

\begin{proposition}[$\Phi^{\FA}_3(\cC)$ on curved CY$_3$]
\label{prop:curved-X}
Let $X$ be a smooth compact CY$_3$. The Stage-1 FA
$\Phi^{\FA}_3(\cC)$ on $X$ is constructed by: (a) covering $X$ by
holomorphic charts $(\cU_\alpha, \Omega_\alpha)$ with flat local
model $(\CC^3, \Omega_0)$; (b) applying the flat-space construction
of Def.~\ref{def:stage1-flat} on each $\cU_\alpha$; (c) gluing via the
CY$_3$-transition functions. The glued FA is canonical (via the
patchwork of contractibly-chosen local data). The global
$\mathbb{E}_3$-structure is defined on the Dolbeault-reduced level
and extends the local $\mathbb{E}_3$-structure of each chart.

Compact-support discipline: the quantisation uses Li--Li 2020
($\emph{Commun.~Math.~Phys.}$ 376, \S5--6), producing a globally
defined effective interaction $I[L]$ on $X$ with local
compactly-supported factorisation algebra of observables.
\end{proposition}

\emph{H5.5 (Deformation theory witness).}

\begin{proposition}[Costello 2013, first-order deformation]
\label{prop:costello-deformations}
The deformation complex of $\hCS(3)$ on $\CC^3$ coupled to $\cC$ is
\[
  \mathrm{Def}(\Phi^{\FA}_3(\cC))
  \;\simeq\;
  \HH^\bullet_{\mathbb{E}_3}(\Phi^{\FA}_3(\cC), \Phi^{\FA}_3(\cC))[3].
\]
For simple Lie $\fg$ on $\CC^3$ with flat volume, first-order
moduli are:
\[
  T_0\cM^{\mathrm{def}} \;=\; H^{0,3}_{\bar\partial,c}(\CC^3)\otimes
  \Sym^2(\fg^\vee)^{\mathrm{inv}}
  \;=\; \CC\cdot\kappa_\fg \text{ (Killing form)},
\]
matching the Yangian deformation moduli at the CY-slice
$\epsilon_1 + \epsilon_2 + \epsilon_3 = 0$.
\end{proposition}

\emph{H5.6--H5.7 (Abelian avatar: CoHA identification).}

\begin{theorem}[Costello--Dimofte--Paquette 2020, reread as Stage-1
identification]
\label{thm:cdp-cohA}
For $\fg = \widehat{\mathfrak{gl}}_1$ (rank-1 abelian), the
all-orders-in-$\hbarb$ BV-quantised $\hCS(3)$ observable FA on
$(\CC^3, \Omega_0)$ is
\[
  \Phi^{\FA}_3(\CoHA(\CC^3))
  \;=\; Y^+(\widehat{\mathfrak{gl}}_1) \;=\; \CoHA(\CC^3)
\]
as an $\mathbb{E}_3$-algebra, where $Y^+$ denotes the positive-nilpotent
part of the Maulik--Okounkov Yangian and the equality is
Schiffmann--Vasserot 2013 (Thm~1.1 loc.~cit.) combined with
Costello--Dimofte--Paquette 2020 (Thm~1.1) identifying the $\hCS$
observables with the Maulik--Okounkov construction.

This pins the Stage-1 output at the canonical CoHA base point and
provides a falsifiable witness: any error in $\Phi^{\FA}_3$ at
$\fg = \widehat{\mathfrak{gl}}_1$ would be visible in the
$Y^+$-product rule (explicit formulas in Schiffmann--Vasserot
\S5.3, loc.~cit.; tested in Vol III compute modules
\texttt{compute/cy\_d\_kappa.py}).
\end{theorem}
\end{heal}

\section*{Synthesis: the five-cycle residue as the Stage-1 theorem}

After five adversarial cycles the Stage-1 construction stabilises to
the following theorem.

\begin{theorem}[Stage 1 of $\Phi_d$, canonical form]
\label{thm:stage1-canonical}
For each $d\in\{2,3,4,5\}$, there exists a functor
\[
  \Phi^{\FA}_d\colon \mathrm{CY}_d\text{-Cat}^{\mathrm{cyc},\mathrm{prop}}
  \longrightarrow \mathbb{E}_d\text{-}\HolFA(X)
\]
of $(\infty,1)$-categories, where the target is the
$(\infty,1)$-category of $\mathbb{E}_d$-holomorphic factorisation
algebras on a CY$_d$-background $X$. The functor is:
\begin{itemize}[leftmargin=2em]
  \item \emph{Defined} by
  $\Phi^{\FA}_d(\cC) := \Obs^{\mathrm{q}}_{\hCS(d)\text{ on }(X,\Omega_X,\fg(\cC))}$,
  i.e.\ the Costello BV-quantised observable FA of
  $(2d)$-dimensional holomorphic Chern--Simons theory coupled to
  $\fg(\cC)$, where $\fg(\cC)$ is the endomorphism $L_\infty$-algebra
  of a compact generator of $\cC$.
  \item \emph{Canonical up to contractible choice} via three
  independently contractible spaces: (i) associator
  $\Phi\in\mathrm{Ass}_d(\QQ)$ (Kontsevich--Tamarkin formality,
  controlled by $\mathrm{GRT}_1(\QQ)$-torsor; Willwacher 2014,
  Fresse--Willwacher 2020); (ii) gauge-fixing operator
  $\bar\partial^*$ (contractible over the space of hermitian
  metrics; Costello 2011 Prop~6.5.3); (iii) scale-family propagator
  $P_L^\varepsilon$ (contractible over $\RR_{>0}$; Costello 2011
  Thm~8.6.1).
  \item \emph{Obeys $\mathbb{E}_d$-structure} via tri-layered
  CGL-locality (Thm~\ref{thm:cgl-locality}) + holomorphic Dunn
  additivity (Prop~\ref{prop:hol-dunn}): locally constant in
  holomorphic direction $+$ locally constant in antiholomorphic
  direction after Dolbeault reduction $=$ $\mathbb{E}_d$.
  \item \emph{Compatible with Stage 2} via composition
  $\Phi_d = \mathrm{Sp}^{\mathrm{ch}}_{\Sigma_{d-1}, C}\circ
  \Phi^{\FA}_d$, where Stage 2 is factorisation homology over
  $\Sigma_{d-1}\subset X$ followed by restriction to a reference
  curve $C\subset X$. Stage 2 is a separate theorem;
  $\Phi_d$ at the functor level is CONJECTURAL, at the value level
  is PROVED for $d\in\{1,2\}$ and K3-fibered Class A at $d = 3$.
  \item At $d = 3$, the classical BV datum is
  $\cA = c + A^{(0,1)} + A^{*(0,2)} + c^{*(0,3)} \in
  \Omega^{0,\bullet}(X,\fg)[1]$, action
  $S_{\mathrm{cl}} = \int_X \Omega_X\wedge \langle\cA,
  \bar\partial\cA + \tfrac{1}{3}[\cA,\cA]\rangle$, $(-1)$-shifted
  symplectic via Serre duality; the Bochner--Martinelli propagator
  generates the FA through Costello--Li 2016 Prop~5.2; the
  $\mathbb{E}_3$-structure is realised by one $\mathbb{E}_1$ per
  holomorphic coordinate after Dolbeault reduction
  (Gwilliam--Williams 2021 Thm~1.1).
\end{itemize}
\end{theorem}

\begin{remark}[Shift-law table, pinned]
\label{rem:shift-table}
The dimensional ladder $(d, \mathrm{AKSZ\ shift}, \mathbb{E}\text{-label})$:
\[
  (2, -2, \mathbb{E}_2), \quad
  (3, -1, \mathbb{E}_1) \text{ on BV} \; / \; \mathbb{E}_3 \text{ on FA}, \quad
  (4, 0, \mathbb{E}_0) \text{ on BV} \; / \; \mathbb{E}_4 \text{ on FA}, \quad
  (5, +1, \mathbb{E}_5\text{-Poisson}) \text{ on BV} \; / \; \mathbb{E}_5 \text{ on FA}.
\]
At $d = 5$ the BV pairing is $+1$-shifted, corresponding to an
$\mathbb{E}_5$-Poisson structure (CPTVV 2017 \S3) on observables; on
the FA side, the $\mathbb{E}_5$-label survives, but the underlying
dgca is Poisson, not symplectic. At $d \geq 6$ the AKSZ shift exceeds
$+1$ and CPTVV shifted-Poisson geometry does not uniformly extend;
Stage 1 is provisional beyond $d = 5$.

The two labels (AKSZ-BV shift on the classical side versus
$\mathbb{E}$-level on the FA side) are \emph{different invariants};
the AKSZ shift measures the symplectic/Poisson degree of
$\omega_{\BV}$, while the $\mathbb{E}$-label measures the
factorisation-algebraic structure on observables. They happen to
coincide at $d = 5$ ($\mathbb{E}_5$ from both sides) but not at
$d < 5$. See Prop.~\ref{prop:shift-fa}.
\end{remark}

\begin{remark}[What Stage 1 is not]
\label{rem:not-stage1}
\begin{itemize}[leftmargin=2em]
  \item Stage 1 is \emph{not} the $\mathbb{E}_1$-chiral algebra on a
        curve; that is Stage 2.
  \item Stage 1 is \emph{not} a numerical invariant ($\kappa$,
        shadow tower, BKM-weight, etc.); those are Vol I / Vol III
        outputs downstream of the full composition $\Phi_d$.
  \item Stage 1 is \emph{not} the Drinfeld double of the chiral
        algebra; $G(\cC) = D(Y^+(\cC))$ is a separate construction in
        Vol III.
  \item Stage 1 does \emph{not} require the CY$_d$-category $\cC$ to
        be geometric (i.e., $D^b\Coh$ of something); it accepts any
        cyclic $A_\infty$-category with Mukai pairing and a
        compact generator.
\end{itemize}
\end{remark}

\begin{remark}[Cross-volume scope]
\label{rem:cross-vol-scope}
The KT formality hypothesis has been adversarially audited in
Vol I cache 434--441 (GRT/KT/Willwacher convention discipline), in
Vol II V2-AP41 (stratified Humbert formality obstruction; smooth KT
does not transport verbatim), and in Vol III AP-CY142 (HH-admissibility
discipline). The CGL locality hypothesis has been adversarially
audited in Vol I AP932 / cache 450 (Dunn real-vs-complex at $\CC^3$),
in Vol II V2-AP139 (Dunn real-vs-complex companion), and in Vol III
AP-CY53 ($\pi_1(\Conf_2)$ ordered-vs-unordered). All three volumes
agree on the tri-layered CGL statement and the
$\mathrm{GRT}_1(\QQ)$-torsor KT statement.
\end{remark}

\section*{Falsification regime}

Stage 1 is a theorem at chain level (Gwilliam--Williams 2021 Thm~1.1)
and at $(\infty,1)$-level (Williams 2017 Thm~4.4.1 + Fresse--Willwacher
2020 Thm~A). It is falsified by the following checks, each
operational:
\begin{itemize}[leftmargin=2em]
  \item \emph{(F1) Bochner--Martinelli normalisation.} The
        $(2\pi i)^{-3}$ prefactor on $P_{\BM}$ on $\CC^3$ is
        determined by the normalisation of the holomorphic volume
        $\Omega_0 = dz_1\wedge dz_2\wedge dz_3$; any change in
        convention propagates to a consistent change in observable
        normalisation. (Cache entry 266, 267 citation discipline.)
  \item \emph{(F2) $\fg = \widehat{\mathfrak{gl}}_1$ Yangian check.}
        $\Phi^{\FA}_3(\CoHA(\CC^3)) = Y^+(\widehat{\mathfrak{gl}}_1)$
        computed by two independent means: direct Feynman diagrams
        (Costello--Dimofte--Paquette 2020 Thm~1.1) and cohomological
        Hall algebra (Schiffmann--Vasserot 2013 Thm~1.1). Any
        mismatch at $\hbarb$-order $\geq 2$ falsifies the Stage-1
        identification.
  \item \emph{(F3) $\HH^\bullet_{\mathbb{E}_3}$ deformation count.}
        First-order deformations of $\Phi^{\FA}_3$ at $\fg$ simple
        compact semisimple: the moduli is $\CC\cdot\kappa_\fg$
        (Killing form, 1-dimensional). Any additional first-order
        moduli direction falsifies the theorem.
  \item \emph{(F4) Dunn-additivity factor count.}
        $\Phi^{\FA}_3(\cC)^{\bar\partial} \simeq
        \Phi^{\FA}_1(\cC)^{\bar\partial, (1)}
        \boxtimes_D \Phi^{\FA}_1(\cC)^{\bar\partial, (2)}
        \boxtimes_D \Phi^{\FA}_1(\cC)^{\bar\partial, (3)}$ (three
        factors, one per holomorphic direction). Any Dunn-factor
        count $\neq 3$ on $\CC^3$ after Dolbeault reduction falsifies.
  \item \emph{(F5) Compactness/non-compactness comparison.} On
        $\CC^3$ (non-compact), the FA uses compactly-supported
        Costello--Gwilliam locality (Prop~8.2.1 loc.~cit.); on
        compact CY$_3$ (quintic, Calabi--Yau hypersurface in weighted
        $\PP^4$), the FA uses Costello 2011 Thm~13.4.1. The two
        outputs agree on the overlap (neighbourhoods of flat patches)
        modulo the global gluing data from the CY-transition
        functions.
\end{itemize}

\section*{First-principles-to-primary-literature dictionary}

The table below maps each step of Stage 1 to the primary paper that
proves it.

\begin{center}
\small
\begin{tabular}{|l|l|l|}
\hline
\textbf{Step} & \textbf{Content} & \textbf{Primary literature} \\
\hline
S1 & $(2d)$-real-dim BV datum of $\hCS(d)$ & Costello--Li 2016 \S2; Costello 2013 \S3 \\
S2 & $(-1)$-shifted symplectic via Serre duality & CPTVV 2017 \S3; Pantev--To\"en--Vaqui\'e--Vezzosi 2013 \\
S3 & AKSZ shift $d - 4$ ladder & CPTVV 2017 Thm~3.1; Calaque 2013 \\
S4 & BV quantisation, effective interaction $I[L]$ & Costello 2011 Thm~13.4.1 \\
S5 & Bochner--Martinelli propagator on $\CC^3$ & Costello--Li 2016 Prop~5.2 \\
S6 & Compactly supported observables FA & CGII Prop~8.2.1; CGII \S10.2 \\
S7 & CGII Thm~5.0.1: FA of observables has $\mathbb{E}_{2d}$ structure & CGII Thm~5.0.1 \\
S8 & Holomorphic translation equivariance, Williams 2017 & Williams 2017 Thm~4.4.1 \\
S9 & Dolbeault reduction, $\mathbb{E}_d$-label & Williams 2017 \S4.4; Gwilliam--Williams 2021 Thm~1.1 \\
S10 & Dunn additivity $\mathbb{E}_d = \mathbb{E}_1^{\otimes_D d}$ & Lurie HA 5.1.2.2; Gwilliam--Williams 2021 \S3 \\
S11 & KT formality of Hochschild $\mathbb{E}_d$-structure & Kontsevich 1999; Tamarkin 2003 Thm~2.1 \\
S12 & $\mathrm{GRT}_1(\QQ)$-torsor of associators & Willwacher 2014 Thm~1.2; Fresse--Willwacher 2020 Thm~A \\
S13 & HKR bridge $\HH^\bullet(D^b\Coh(X)) \simeq H^\bullet(X, \wedge^\bullet TX)$ & HKR 1962; Swan 1996; Caldararu 2005 \\
S14 & Deformation complex = $\HH^\bullet_{\mathbb{E}_3}[3]$ & Costello 2013 Thm~2.1 \\
S15 & $\fg = \widehat{\mathfrak{gl}}_1$ all-orders = $Y^+$ & Costello--Dimofte--Paquette 2020 Thm~1.1; Schiffmann--Vasserot 2013 \\
S16 & Canonicity up to contractible choice, $(\infty,1)$-cat level & Lurie HA; Fresse--Willwacher 2020 \\
\hline
\end{tabular}
\end{center}

Every row is a theorem in the cited primary paper; Stage 1 is
obtained by composing these rows in the order S1$\to$S16.

\section*{Concordance check}

Stage 1 of $\Phi_d$ as stated here is consistent with:
\begin{itemize}[leftmargin=2em]
  \item \texttt{CLAUDE.md} \S ``Programme identity'': Vol III
        two-stage factorisation
        $\Phi_d = \mathrm{Sp}^{\mathrm{ch}}_{\Sigma_{d-1},C}
        \circ \Phi^{\FA}_d$ (this paper constructs $\Phi^{\FA}_d$).
  \item \texttt{memory/platonic\_ideal\_reconstituted\_2026\_04\_17.md}
        \S3 G4, \S5.5 shadow functor: the diagram in
        $(\infty,2)$-Cat whose top 2-cell is
        $\Phi_d \simeq \mathrm{Sp}^{\mathrm{ch}}\circ\Phi^{\FA}_d$
        (this paper inscribes the top row of that diagram).
  \item \texttt{calabi-yau-quantum-groups/notes/platonic\_synthesis\_waves\_11\_through\_16.tex}
        Thm~\texttt{wn:thm:plat-hCS-classical}--\texttt{wn:thm:plat-dualizability}
        (this paper reads those theorems as the $\hCS(3)$ avatar of
        Stage 1 at $d = 3$).
  \item \texttt{notes/6d\_hCS\_audit\_2026\_04\_22\_cross\_volume\_errors.md}:
        the cross-volume error synthesis catalogues 24 specific errors
        in earlier drafts (quartic-for-quadratic Casimir, $\phi_{-2,1}$
        polar coefficient, etc.); this paper avoids all of them by
        citing only compactly-supported CGL-locality on $\CC^3$ and
        deferring numerical claims on $K3\times E$ to Vol III.
  \item \texttt{notes/wave17\_opus\_20260424/opus\_13\_6d\_hCS\_CFG\_machine.tex}:
        the companion deep derivation of the 6D hCS CFG machine on
        $\CC^3$ (Bochner--Martinelli explicit, Deligne-exceptional
        anomaly, $E_3$-Koszul dualisability). This paper restricts to
        the Stage-1-as-functor viewpoint and does not duplicate the
        CFG machine's internal derivations.
\end{itemize}

\section*{Bibliography anchors (primary)}

\begin{enumerate}[leftmargin=2em]
  \item Calaque, Pantev, To\"en, Vaqui\'e, Vezzosi, ``Shifted Poisson
        structures and deformation quantization,''
        \emph{J.~Topol.} 10 (2017), arXiv:1506.03699. (Shift law,
        CPTVV.)
  \item Costello, \emph{Renormalization and Effective Field Theory},
        AMS Mathematical Surveys and Monographs 170, 2011. (BV
        quantisation machinery, Thm~13.4.1.)
  \item Costello, ``Supersymmetric gauge theory and the Yangian,''
        arXiv:1303.2632, 2013. (Deformation of 6D hCS.)
  \item Costello, Dimofte, Paquette, ``Factorization algebras in the
        cohomological Hall algebra,'' \emph{Commun.~Math.~Phys.} 385
        (2021), arXiv:2007.14610. (CoHA identification at
        $\widehat{\mathfrak{gl}}_1$.)
  \item Costello, Gwilliam, \emph{Factorization Algebras in Quantum
        Field Theory}, Vols I \& II, Cambridge UP, 2017, 2021.
        (Locality theorems, $\mathbb{E}_n$-structure on observables.)
  \item Costello, Li, ``Twisted supergravity and its quantization,''
        arXiv:1605.09930, 2016. (Bochner--Martinelli propagator,
        Prop~5.2.)
  \item Fresse, Willwacher, ``The intrinsic formality of
        $E_n$-operads,''
        \emph{J.~Eur.~Math.~Soc.} 23 (2021), arXiv:1503.08699, and
        sequel \emph{Comp.~Math.} 156 (2020). (GRT torsor of
        $\mathbb{E}_d$-formality.)
  \item Gwilliam, Williams, ``Higher Kac--Moody algebras and
        symmetries of holomorphic field theories,''
        \emph{Adv.~Theor.~Math.~Phys.} 25 (2021) 129--239,
        arXiv:1810.06363. (Explicit $\mathbb{E}_3$-algebra on
        $\CC^3$.)
  \item Hochschild, Kostant, Rosenberg, ``Differential forms on
        regular affine algebras,''
        \emph{Trans.~AMS} 102 (1962), 383--408. (HKR isomorphism.)
  \item Kontsevich, ``Deformation quantization of Poisson manifolds,''
        \emph{Lett.~Math.~Phys.} 66 (2003), arXiv:q-alg/9709040,
        1999. ($\mathbb{E}_2$-formality for smooth functions.)
  \item Li, ``Vertex algebras and quantum master equation,''
        \emph{J.~Differential Geom.} 114 (2020). (Feynman-diagram
        regularisation.)
  \item Li, Li, ``Twisted supergravity revisited: a partial
        perspective on its quantization,''
        \emph{Commun.~Math.~Phys.} 376 (2020), arXiv:1809.11014.
        (RG-flow analysis of 6D hCS.)
  \item Lurie, \emph{Higher Algebra}, online notes, 2017. (HA 5.1.2.2
        Dunn additivity; HA 5.3.4.10 locally constant FA $=$
        $\mathbb{E}_n$-algebra.)
  \item Maulik, Okounkov, ``Quantum groups and quantum cohomology,''
        arXiv:1211.1287, 2012. (Yangian framework.)
  \item Pantev, To\"en, Vaqui\'e, Vezzosi, ``Shifted symplectic
        structures,''
        \emph{Publ.~IHES} 117 (2013). (Shifted symplectic geometry.)
  \item Schiffmann, Vasserot, ``The elliptic Hall algebra,
        cherednik Hecke algebras and Macdonald polynomials,''
        \emph{Compositio Math.} 147 (2011); ``The elliptic Hall
        algebra and the $K$-theory of the Hilbert scheme of $\CC^2$,''
        \emph{Duke Math.~J.} 162 (2013); ``Cohomological Hall algebras,
        vertex algebras, and instantons,''
        arXiv:1705.01435, 2017. (CoHA.)
  \item Tamarkin, ``Action of the Grothendieck--Teichm\"uller group
        on the Hochschild complex,''
        preprint 2003, arXiv:math/0301098. (KT $\mathbb{E}_d$-formality.)
  \item Voronov, ``Notes on universal algebra,''
        \emph{Ann.~Sci.~ENS} (2000). (Gerstenhaber structure on
        Hochschild.)
  \item Williams, \emph{The Holomorphic $\sigma$-Model and its
        Symmetries}, PhD thesis, Northwestern University, 2017.
        (Thm~4.4.1 holomorphic Dunn.)
  \item Williams, ``Renormalization for holomorphic field theories,''
        \emph{Commun.~Math.~Phys.} 380 (2020), arXiv:1809.02661.
        (Factorisation-homology path.)
  \item Willwacher, ``M.~Kontsevich's graph complex and the
        Grothendieck--Teichm\"uller Lie algebra,''
        \emph{Invent.~Math.} 200 (2014). (GRT torsor of associators.)
\end{enumerate}

\end{document}
